\chapter{Khởi Động Sau Giờ Kiểm Tra}
\chapterintro{Giải tỏa vừa đủ rồi đi tiếp}{Sau kiểm tra, học sinh thường còn nguyên cảm xúc. Khởi động lúc này phải giúp chuyển trạng thái chứ không kéo dài chuyện cũ.}

\section{Một từ xả áp}
\begin{khoidongbox}{Hoạt động khởi động}
Mời mỗi em nói đúng một từ cho trạng thái hiện tại sau bài kiểm tra.
\end{khoidongbox}
\begin{visaohieuqua}
Xả một lượng nhỏ giúp lớp nhẹ đầu hơn. Một từ là đủ để nói thật mà không kéo tiết học chậm lại.
\end{visaohieuqua}
\begin{cachlam5phut}
Lấy 5 đến 7 từ, cười nhẹ nếu cần, rồi nói rõ: \textit{``Rồi, mình sang phần mới.''}
\end{cachlam5phut}
\ngancach

\section{Tách kiểm tra khỏi bài mới}
\begin{khoidongbox}{Hoạt động khởi động}
Nói dứt khoát: \textit{``Phần kiểm tra dừng ở đây, phần học mới bắt đầu từ đây.''}
\end{khoidongbox}
\begin{visaohieuqua}
Học sinh cần được dẫn qua ranh giới tâm lý. Câu tách nhịp rõ ràng giúp các em buông bớt chuyện cũ.
\end{visaohieuqua}
\begin{cachlam5phut}
Sau câu này, cho lớp làm ngay một nhiệm vụ cực dễ để lấy lại cảm giác vào bài được.
\end{cachlam5phut}
\ngancach

\section{Một câu thắng nhỏ đầu tiên}
\begin{khoidongbox}{Hoạt động khởi động}
Mở phần mới bằng một câu hỏi rất dễ hoặc một lựa chọn chắc thắng cho cả lớp.
\end{khoidongbox}
\begin{visaohieuqua}
Sau kiểm tra, học sinh rất cần cảm giác mình vẫn làm được. Một chiến thắng nhỏ giúp lớp hồi năng lượng nhanh hơn.
\end{visaohieuqua}
\begin{cachlam5phut}
Chỉ cần 1 câu dễ, sau đó mới tăng dần độ khó. Không mở ngay bằng phần nặng nhất.
\end{cachlam5phut}
\ngancach

\section{Chọn bài mới theo kiểu nhẹ hay nhanh}
\begin{khoidongbox}{Hoạt động khởi động}
Hỏi lớp: \textit{``Sau bài kiểm tra, mình muốn vào phần mới theo kiểu nhẹ hơn hay nhanh hơn?''}
\end{khoidongbox}
\begin{visaohieuqua}
Câu hỏi trao một phần lựa chọn cho học sinh. Khi được chọn nhịp vào bài, các em hợp tác hơn rõ rệt.
\end{visaohieuqua}
\begin{cachlam5phut}
Cho lớp giơ tay chọn, rồi điều chỉnh tốc độ mở bài theo đa số nhưng vẫn giữ mục tiêu tiết học.
\end{cachlam5phut}
\ngancach

\section{Đặt tên cho tâm trạng sau kiểm tra}
\begin{khoidongbox}{Hoạt động khởi động}
Mời lớp đặt một cái tên vui cho không khí sau bài kiểm tra vừa xong.
\end{khoidongbox}
\begin{visaohieuqua}
Đặt tên cho cảm xúc giúp học sinh nhìn nó nhẹ hơn. Lớp thường bật cười và dễ đi tiếp hơn nhiều.
\end{visaohieuqua}
\begin{cachlam5phut}
Lấy 2 đến 3 cái tên, chọn một tên hay nhất, rồi chốt: \textit{``Rồi, mình ra khỏi bộ phim đó và vào bài mới.''}
\end{cachlam5phut}